% ===========================================================================
%  E10 — Painel da Unooria
% ===========================================================================
\chapter{E10 \textperiodcentered{} Painel da Unooria}
\label{ep:e10}

\begin{fichaepico}
\begin{campos}
  \item[Objetivo] A equipe da Unooria decide o que existe em cada rede, qual
    versão de cada missão está no ar e quanto de IA cada perfil pode usar ---
    e recebe os interessados do site.
  \item[Para quem] Administrador Unooria. O Administrador do contrato só lê o
    histórico.
  \item[Valor] É o que permite vender produtos separadamente, desligar uma
    missão com problema em minutos e voltar uma versão sem esperar o time.
  \item[Estado] \pronto{} --- faltam desligar por turma ou por prazo e alguns
    avisos.
\end{campos}
\end{fichaepico}

\begin{tabularx}{\linewidth}{@{}L{16mm}X L{20mm}@{}}
\toprule
\headrow \thd{ID} & \thd{Feature} & \thd{Estado}\\
\midrule
\id{F10.1} & Ligar e desligar produtos e missões & \pronto\\
\zebra \id{F10.2} & Histórico que a rede também lê & \pronto\\
\id{F10.3} & Versões das missões & \pronto\\
\zebra \id{F10.4} & Uso de IA & \pronto\\
\id{F10.5} & Prospecção & \pronto\\
\zebra \id{F10.6} & Experiências arquivadas & \pronto\\
\bottomrule
\end{tabularx}

\begin{nota}
\textbf{Por que só a Unooria opera.} O Administrador do contrato pode quase
tudo na rede dele --- mas não pode ligar o que a rede não contratou. É a única
capacidade que a Unooria tem e a rede não.
\end{nota}

% ---------------------------------------------------------------------------
\feature{F10.1}{Ligar e desligar produtos e missões}{\pronto}

\begin{historia}{H10.1.1}{Desligar um produto ou uma missão}{\pronto}
\hu{Administrador Unooria}{desligar um produto ou uma missão numa rede, numa
regional ou numa escola}{respeitar o que foi contratado ou tirar do ar uma
missão com problema}
\aceite
\begin{itemize}
  \item Em \enquote{Recursos}, escolho o Escopo: \enquote{Rede (contrato
    inteiro)}, uma regional ou uma escola.
  \item Vejo os cinco produtos e, dentro de cada um, as missões (o agente de IA
    do professor não tem missões dentro). Cada
    interruptor tem três estados: \enquote{Ligado}, \enquote{Desligado} e
    \enquote{Herda de cima} (na rede, \enquote{Padrão (ligado)}).
  \item Por padrão tudo está ligado; o escopo mais específico vence (escola,
    depois regional, depois rede).
  \item O produto é um teto: desligado, as missões dentro dele ficam
    desligadas, com o motivo ao lado.
  \item Desligar exige motivo; na rede, o aviso diz \enquote{Isto vale para a
    REDE INTEIRA}.
  \item A missão desligada sai da grade da turma (não fica com cadeado), e
    para a criança a mensagem é a mesma de missão ainda não aberta.
  \item Religar devolve a sala como estava: as liberações do professor nunca
    são apagadas.
  \item Com o produto desligado, cada pessoa continua podendo apagar os
    próprios dados.
\end{itemize}
\end{historia}

\begin{historia}{H10.1.2}{Missão nova já nasce com interruptor}{\pronto}
\hu{Administrador Unooria}{que toda missão ou ferramenta nova apareça sozinha
no painel}{nunca ter uma missão no ar que eu não consiga desligar}
\aceite
\begin{itemize}
  \item A missão nova aparece dentro do produto a que pertence, sem ninguém
    precisar cadastrá-la à mão.
  \item O teste de aceite é desligá-la e ver que ela some da sala e do
    estudante.
\end{itemize}
\end{historia}

\begin{historia}{H10.1.3}{Desligar por turma ou por prazo}{\afazer}
\hu{Administrador Unooria}{desligar um recurso só para uma turma, ou até uma
data}{atender pilotos e contratos temporários sem mexer na escola inteira}
\aceiteproposta
\begin{itemize}
  \item O Escopo passa a incluir turma.
  \item Ao desligar, posso informar até quando; na data, o interruptor volta a
    herdar de cima e isso fica no histórico.
\end{itemize}
\end{historia}

\begin{historia}{H10.1.4}{Desligar o agente de IA do professor numa rede}{\parcial}
\hu{Administrador Unooria}{ligar e desligar o agente de IA do professor como
qualquer produto}{oferecê-lo só a quem contratou, ou tirá-lo do ar num
problema}
\aceite
\begin{itemize}
  \item O agente aparece em \enquote{Recursos} como um produto, com o mesmo
    interruptor de três estados e o mesmo escopo (rede, regional, escola).
  \item Desligado, o botão do agente some das telas do professor naquela
    escola, e nenhuma conversa nova nem mensagem passa.
  \item Mesmo desligado, o professor continua podendo apagar as próprias
    conversas.
\end{itemize}
\pendencia{o interruptor já funciona; o agente em si ainda não está no ar nos
ambientes da rede (\id{F08.4}).}
\end{historia}

% ---------------------------------------------------------------------------
\feature{F10.2}{Histórico que a rede também lê}{\pronto}

\begin{historia}{H10.2.1}{Saber quem mudou o quê, e por quê}{\pronto}
\hu{Administrador do contrato}{ver tudo o que a Unooria ligou, desligou ou
mudou na minha rede}{explicar à secretaria por que algo sumiu}
\aceite
\begin{itemize}
  \item O histórico registra quem, o quê, quando e o motivo; o nome de quem
    mudou fica como era no dia.
  \item Ninguém altera o histórico.
  \item O Administrador do contrato lê em \enquote{Recursos da rede}, só
    leitura: \enquote{Quem opera estes interruptores é a Unooria; esta tela é
    o registro.}
  \item As coordenações não veem esta tela.
\end{itemize}
\end{historia}

% ---------------------------------------------------------------------------
\feature{F10.3}{Versões das missões}{\pronto}

\begin{historia}{H10.3.1}{Voltar para a versão anterior}{\pronto}
\hu{Administrador Unooria}{voltar uma missão para a versão anterior}{resolver
em minutos um problema que apareceu numa versão nova}
\aceite
\begin{itemize}
  \item A janela Versão mostra as versões publicadas, com a atual marcada
    \enquote{(no ar)}, e \enquote{Ativar para a rede}.
  \item Só se volta para versões que já existiram; a volta restaura também as
    configurações daquela versão.
  \item A tela diz se conferiu que a versão está de fato publicada.
  \item Enquanto isso, a missão aparece com o aviso \enquote{Esta experiência
    está em rollback}.
\end{itemize}
\end{historia}

\begin{historia}{H10.3.2}{Testar uma versão nova numa escola antes}{\pronto}
\hu{Administrador Unooria}{fixar uma versão numa escola ou regional}{testar
uma novidade num piloto antes de levar para a rede}
\aceite
\begin{itemize}
  \item Fora do escopo da rede: \enquote{Fixar neste escopo} e \enquote{Soltar
    a fixação}.
\end{itemize}
\end{historia}

\begin{historia}{H10.3.3}{Ser lembrado de uma versão fixada há muito tempo}{\afazer}
\hu{Administrador Unooria}{ser avisado quando uma versão fixada ficar muito
tempo sem ser solta}{não deixar uma escola esquecida numa versão antiga}
\aceiteproposta
\begin{itemize}
  \item O painel destaca fixações com mais de 30 dias.
\end{itemize}
\end{historia}

% ---------------------------------------------------------------------------
\feature{F10.4}{Uso de IA}{\pronto}

\begin{historia}{H10.4.1}{Definir quanto de IA cada perfil pode usar}{\pronto}
\hu{Administrador Unooria}{definir o limite diário de IA dos estudantes e dos
professores}{manter o custo previsível sem travar a aula}
\aceite
\begin{itemize}
  \item Dois cartões, \enquote{Estudantes} e \enquote{Professores}, com o
    limite em \enquote{Tokens por dia} traduzido em uso real (partidas da
    Missão Tempestade, planos anuais).
  \item Valores iniciais: cerca de 3 partidas por dia para o estudante e 6 a 9
    planos anuais para o professor.
  \item Mudar exige motivo e avisa \enquote{Vale para a rede inteira}; vale na
    próxima chamada, sem esperar o time.
  \item Zero é recusado: \enquote{para fechar a IA de um produto, desligue o
    produto em Recursos}.
  \item Toda mudança vai para o histórico. A rede não vê nem altera.
\end{itemize}
\end{historia}

% ---------------------------------------------------------------------------
\feature{F10.5}{Prospecção}{\pronto}

\begin{historia}{H10.5.1}{Receber os interessados do site}{\pronto}
\hu{Administrador Unooria}{ver quem se cadastrou no site}{entrar em contato
com as redes interessadas}
\aceite
\begin{itemize}
  \item \enquote{Prospecção} lista Quando, Nome, E-mail, Telefone,
    \enquote{Rede, secretaria ou escola} e Cargo, dos mais recentes aos mais
    antigos (até 5 mil).
  \item \enquote{Atualizar} e \enquote{Baixar CSV}.
  \item Só a Unooria vê esta lista.
\end{itemize}
\end{historia}

\begin{historia}{H10.5.2}{Eliminar um cadastro a pedido}{\afazer}
\hu{Administrador Unooria}{apagar o cadastro de quem pedir}{atender o direito
de eliminação da LGPD sem depender do time técnico}
\aceiteproposta
\begin{itemize}
  \item Cada linha tem \enquote{Eliminar a pedido do titular}, com
    confirmação; o cadastro nunca é editado, só eliminado.
\end{itemize}
\pendencia{hoje o pedido é atendido por um caminho manual.}
\end{historia}

% ---------------------------------------------------------------------------
\feature{F10.6}{Experiências arquivadas}{\pronto}

Quando uma missão nova ocupa o lugar de outra na trilha, a antiga não é apagada:
ela é \textbf{arquivada}. Sai da sala, do catálogo e do painel de recursos, e os
dados que as turmas produziram ficam guardados.

\begin{historia}{H10.6.1}{Ver o que saiu do produto sem ser apagado}{\pronto}
\hu{Administrador Unooria}{ver quais missões foram arquivadas e o que ficou
guardado delas}{saber o que a plataforma já teve e responder por aqueles dados}
\aceite
\begin{itemize}
  \item \enquote{Arquivados} lista cada missão arquivada com o nome, o código de
    origem, onde ficava na trilha, a última versão, os dados guardados (turmas e
    estudantes), o motivo e a data.
  \item Os números de dados guardados são os do dia do arquivamento, e a tela
    diz isso.
  \item A lista só mostra: ninguém joga uma missão arquivada, e trazê-la de
    volta é trabalho do time técnico.
  \item Só a Unooria vê esta lista.
\end{itemize}
\end{historia}
